% =====================================================================
\chapter{结论}\label{ch:conclusion}
% =====================================================================

\section{四个问题的回答}

\begin{problem}[底层确定性的可推广性]
结论分三层。底层之所以近乎唯一，是因为它被\emph{过约束}：约束数不少于设计自由度，
解集退化为离散点集（命题 \ref{prop:overconstrain}）。高层之所以多解，是因为问题
层面的约束远少于自由度，解集是一个高维流形。但高层仍然收敛，因为存在\emph{第二组
约束}（来自制造与维护该系统的人类组织）把流形压缩成少数几个尖峰
（原则 \ref{pr:org}）。

因此准确的回答是：存在收敛，但不是朝向工程最优，而是朝向\emph{有界组织可导航的
结构}。收敛的载体是\emph{关系}而非组件（原则 \ref{pr:relations}）；收敛的强度关于
领域绑定度单调递减（猜想 \ref{conj:strength}）。
\end{problem}

\begin{problem}[规模的强制性]
结论是一个\emph{条件必然性}结果。决定是否需要统一基座的自变量是活跃实体数与共享
资源数的乘积，而不是代码行数（定理 \ref{th:integration}、推论 \ref{cor:threshold}）。
而六项职能（身份与生命周期、分配与回收、调度、命名与寻址、通信与同步、故障隔离与
恢复）是操作语义良定义的必要条件（定理 \ref{th:six}）。它们可以被自建，也可以被
外借；判据是其所承载的不变量是否是领域关键的（推论 \ref{cor:buildadopt}）。
\end{problem}

\begin{problem}[智能的理解与设计]
结论是：理解是逆向压缩，设计是正向压缩，二者是同一结构的两个读取方向。理解的最优
下潜深度由率失真驻点条件给出（命题 \ref{prop:stopping}）；设计的可测量目标是读集
期望（定义 \ref{def:readset}），其最小化要求不变量的单点强制（定理
\ref{th:centralise}）。抽象是否泄漏是可测的，而不只是可感的（命题
\ref{prop:leakpair}、推论 \ref{cor:diagnose}）。对 AI 而言，优势不在于读得更多，
而在于探测得更快（命题 \ref{prop:active}）。
\end{problem}

\begin{problem}[无人审阅条件下的正确性]
结论是：人不再逐行审阅实现是可达的，人不必再为目标的表达负责是不可达的。
"正确"必须被拆成五个维度（定义 \ref{def:five-correctness}）；中间过程必须被定义为
七段可验证细化链（定义 \ref{def:chain}）；信任必须是两种生成模式的复合
（定义 \ref{def:two-layer}）；制品绑定是不可省略的第四支柱（式
\eqref{eq:binding}）。综合保证条件式（命题 \ref{prop:assurance}）给出了完整的回答。
\end{problem}

\section{压缩作为主线}

把全文收拢，得到一条单一的主线：\term{压缩}。

\begin{itemize}
\item 复杂系统用\emph{层级}压缩自己，于是每个大系统都长出一个域内内核
      （第 \ref{ch:infrastructure} 章）；
\item 有限智能用\emph{理解}压缩信息，于是有限的认知装得下无限的世界
      （第 \ref{ch:comprehension} 章）；
\item 工程用\emph{架构}压缩变更成本，于是局部的修改不需要全局的理解
      （第 \ref{ch:design} 章）；
\item 文明用\emph{规则}压缩信任，于是来源不明的百万行代码可以被安全运行
      （第 \ref{ch:intent} 章）。
\end{itemize}

而压缩的机制在四种情形下是同一个：按抽象映射的核做商
（命题 \ref{prop:oblivious}）。\emph{一个层级暴露给上层的全部信息，就是下层行为
在这个核下的商。}

\section{核心命题}

如果要把全文压缩为一句可被检验的话，本文选择这一句：

\begin{principle}[核心命题]\label{pr:core}
一个复杂系统能被有限智能理解，当且仅当它可被压缩；而它可被压缩，当且仅当它存在
一套尺度分离，使得每一层都拥有一个明确的、被声明的\emph{无视集}。
\end{principle}

这句话有三个直接后果。

第一，它把"可理解性"从一个主观感受变成了一个结构性质：可理解性等价于存在一个
商映射，在该映射下任务相关的预测能力被保留。这使可理解性成为设计的\emph{目标函数}
而非事后评价。

第二，它把"架构"重新定义为\emph{人为制造的尺度分离}。自然界的分层由尺度分离与
重整化免费获得；软件的分层不会自然涌现，必须由接口、封装与不变量付费制造，并且
随时会漏。因此架构的核心交付物不是模块列表，而是\emph{每一层的无视集声明}。

第三，它对 AI 给出一个具体的行动准则：不要试图穷尽底层信息量，而应维护一个分层的、
可证伪的、能被反例更新的压缩模型；用探测而非阅读获取梯度；以满足任务充分性作为
停止条件。

\section{结语}

全文始于三个工程直觉，落在一个关于结构与保证的联合命题上。确定的不是答案，而是
\emph{哪些问题没有答案}，以及为什么：

\begin{itemize}
\item 意图与规格之间的鸿沟没有机械解，其根据是认识论的而不是逻辑的；它实质上是
      框架问题，因而只在交互中可回避；
\item 机械保证的信任基础可以被缩小，但不能被消除；且"验证便宜"必须分解为检查时间
      与证书尺寸两个独立命题，后者才是常见瓶颈；
\item 共识保证的可靠性受相关错误模式支配，因此\emph{多样性}而不是\emph{数量}
      才是设计变量，而多样性最终是激励设计问题，不是统计问题；
\item 大系统的结构收敛是概率性的、分领域的、受组织约束塑形的，而不是唯一的；
\item 架构的职责不是降低单个变更的成本，而是降低该成本随规模增长的\emph{指数}。
\end{itemize}

把这几条并列起来看，共同的形状是：\emph{复杂度压缩是规模驱动的必然，压缩的形态是
领域驱动的偶然；而保证的边界，恰好落在"能够被机器判定的"与"必须由承担后果的主体
确认的"之间。}

代价是放弃一个可以用一条定理表达的图景；收益是每一个主张都落在它能被检验的位置上。
对一个关于"如何在不可信的复杂性之上建立可信结论"的框架，这是唯一自洽的结局：
\emph{它必须通过自己的检验。}
